\section{Interpret temporary values at the right boundary}\label{rgb-working-state}
The public arguments become working pointers and remaining count. A scalar pack
uses \texttt{t0}, \texttt{t1} and \texttt{t2} for channel contributions. Expansion
keeps the numerical packed value in \texttt{t0} and produces red, green and blue in
\texttt{t1}, \texttt{t2} and \texttt{t3}. These are internal roles. On a completed
call, only the specified public result and preserved state are required to match.

Every named region begins as a branch and \texttt{nop}. Replace that placeholder
instead of appending after it. The initial vector gate placeholders choose the
scalar tails, allowing the scalar loops to be completed first. Other placeholders
choose their function's return. Some stages therefore return a correct count with
unfinished output. Use only the explicitly stated construction observation for
such a stage; final grading requires the complete output and ownership contract.

The verification fixtures call the public function and pause before its supplied
return continuation. No caller instruction has overwritten temporary registers or
vector lanes. The provided leaf return consists only of \texttt{jr} and its
\texttt{nop} slot. Inspecting a vector's byte or halfword lanes changes the view,
not its stored bits. The programmer must know which layout the preceding operation
established before interpreting that view.

The code keeps a \texttt{nop} after each retained delayed branch or jump. Preserve
those slots while implementing the conversions. Moving an unrelated instruction
there can change both branch outcomes and invalidate an otherwise correct pixel
calculation. A full pointer remains full width even when the next displacement is
only two, three, sixteen or twenty-four bytes.
